
\subsection{Key observations}

Given a \emph{restricted} \laVASS $V$ in dimension $d$, assuming that
$L(V)\neq \emptyset$, there is a run $\pi$ from an initial
configuration $c$ to a final configuration $c'$. With no loss of
generality, throughout this section, we assume that $\val(c',x_i) \ge
\val(c',x_{i+1})>0$ for all $1\le i<d$. In particular, this implies
that every counter gets incremented at least once along a path
witnessing non-emptiness.

Our first observation is that if along $\pi$ a counter $y_i$ achieves
the first time a non-zero value by taking a $\wTTPO$ labeled
transition, the length of the remaining segment of $\pi$ is bounded by
$O(\log(m_i+1))$, where $m_i$ is the value of counter $x_i$ before the
transition is taken. The reason is that, once $y_i$ has non-zero
value, its value at least doubles when a transition is taken. Hence if
$\pi$ is ``long'' then along $\pi$ there will be loops incrementing a
counter $x_i$ before the corresponding $y_i$ achieves non-zero value.

In the latter scenario, we may actually, subject to some bookkeeping,
discard concrete values of $x_i$ and $y_i$ and only store their
residue classes modulo $\ell_i$, where $\ell_i$ is the length of the
first loop incrementing $x_i$ along $\pi$. In particular, if we are
given a non-accepting run $\pi'$ such that $\val(\pi',x_i) \equiv
\val(\pi',y_i) \bmod \ell_i$ and $\val(\pi',x_i)<\val(\pi',y_i)$ then
$\pi'$ can be turned into a run $\pi''$ where $\val(\pi'',x_i) =
\val(\pi'',y_i)$ by iterating the loop of length $\ell_i$.

There are, however, some further subtleties that need to be taken care
of. Consider the segment $\pi'$ of $\pi$ between the first transition
labeled by $\wPP$ on $x_i$ and the first transition labeled by $\wPP$ on
$x_{i+1}$. If $\pi'$ contains no loop then we are in a situation where
the first loop incrementing $x_i$ is also the first loop incrementing
$x_{i+1}$. This means that the values of $x_i$ and $x_{i+1}$ get
paired together, and hence, for an accepting run, also the values of
$y_i$ and $y_{i+1}$ are paired together. In our approach, we deal with
such circumstances by introducing so-called \emph{$y$-constraints}. A
$y$-constraint of the form $y_i-y_{i+1}=\delta_i$ for some constant
$\delta_i \in \N$ asserts that the counters $y_{i+1}$ and $y_i$ must
eventually have constant difference $\delta_i$ along a run.

Otherwise, if $\pi'$ above contains a loop, the difference between the
values of $x_i$ and $x_{i+1}$ is not necessarily constant, but
lower-bounded by the length $\delta_i$ of the loop-free sub path of
$\pi'$. Thus, in an accepting run, the difference between $y_i$ and
$y_{i+1}$ must also be at least $\delta_i$, which is asserted by a
$y$-constraint of the form $y_i-y_{i+1} \ge \delta_i$.



%% \subsection{Shortest accepting paths can be very long}

%% We present an example of a restricted \laVASS that has an accepting
%% path, but the length of the shortest such path is non-elementary. All
%% the unspecified transitions labels are $\timestwo$ for the y-counters
%% and $id$ for the x-counters with value zero, and $\wPP$ for the counters
%% with value greater than zero, respectively.

%% \begin{figure}[!t]

%% \begin{center}
%% \begin{tikzpicture}[->,shorten >=1pt,auto,node distance=2.2cm,semithick]

%% \tikzstyle{every state}=[draw=black,text=black, minimum size = 1mm]

%% \node[initial by arrow,state,initial text=] (1) {};
%% \node[state] (2) [right of=1] {};

%% \node[state] (4) at (0,-1.1) {};
%% \node[state] (3) [right of=4] {}; 
%% \node[state] (5) [right of=3] {};
%% \node[state, accepting] (9) [right of=5] {};

%% \path 
%%     (1) edge node [align=center]{$y_d:\timestwob$}  (2)
%%     (1) edge [loop above] node {$x_d:\wPP$}(1)
    
%%     (2) edge [loop above] node {$x_{d-1}:\wPP$} (2)
    
%%     (2) edge [dotted] node {} (4)


%%     (4) edge node [align=center,below]{$y_{2}:\timestwob$}  (3)
%%     (4) edge [in=120,out=60,loop below] node {$x_{2}:\wPP$}(4)
    
%%     (3) edge node [align=center,below]{$y_1:\timestwob$}  (5)
%%     (3) edge [in=120,out=60,loop below] node {$x_1:\wPP$}(3)
    
    
%%     (5) edge node [align=center,below]{$y_1:\timestwob$}  (9);
    

    
%% \end{tikzpicture}
%% \end{center}

%% \caption{An \laVASS with an accepting path requiring non-elementary
%%   length.}
%% \label{fig:non-elem}
%% \end{figure}

%% Let $c=(q,m_1,n_1,\dots,m_d,n_d)$ be the final configuration we reach after we take the shortest accepting path of the \laVASS in \Cref{fig:non-elem}. Clearly, $n_1$ (the value of the counter $y_1$) is 3 and so, $m_1$ must be 3 since $c$ is a final configuration. 

%% Now assume that $m_i,n_i\geq2\uparrow\uparrow i$ (here, $\uparrow$ is the Knuth up-arrow notation), for all $i\in\{1,\dots,k-1\}$. Let us consider the values of $m_{k},n_k$:

%% \begin{figure}[h]

%% \begin{tikzpicture}[->,shorten >=1pt,auto,node distance=2.2cm,semithick]

%% \tikzstyle{every state}=[draw=black,text=black, minimum size = 3mm]

%% \node (1) at (2.5,0) {.};
%% \node (2) at (3,0) {.};
%% \node (3) at (3.5,0) {.};
%% \node[state] (4) at (4.5,0) {};
%% \node[state] (5) [right of=4] {}; 
%% \node (6) at (8,0) {.};
%% \node (7) at (8.5,0) {.};
%% \node (8) at (9,0) {.};

%% \path 

%%     (4) edge node [align=center]{$y_{k}:\timestwob$}  (5)
%%     (4) edge [in=120,out=60,loop above] node {$x_{k}:\wPP$}(4)
    
%%     (5) edge [in=120,out=60,loop above] node {$x_{k-1}:\wPP$}(5);
    
    
%% \end{tikzpicture}
%% \centering

%% \label{fig:induction step}
%% \end{figure}

%% Since we know that $m_{k-1}\geq 2\uparrow\uparrow (k-1)$ it means that $n_k\geq 2^{2\uparrow\uparrow (k-1)}=2\uparrow\uparrow k$. However, we know that $m_k=n_k$, so $m_k\geq 2\uparrow\uparrow k$. The key observation is that $x_{k-1}$ needs to take many self-loops to recover the difference to the $y_{k-1}$ counter. However, every time $x_{k-1}$ takes a self-loop, the value of $y_{k}$ doubles and $x_k$ needs to take an exponential number of self-loops in order to recover that, which in turn doubles the value of counter $y_{k+1}$ each time. 

%% We can conclude that $m_d\geq 2\uparrow\uparrow d$ and since $m_d$ is
%% equal to the length of the shortest accepting path (we cannot take any
%% transition in the beginning without updating the value of $m_1$), we
%% have shown that the length of the shortest accepting path of a \laVASS
%% can be non-elementary with respect to the size of the input restricted
%% \laVASS.
